% 第8章 8.1.3：固定宽度 sigma=0.15 的高斯核截面 K(x,c)=exp(-(x-c)^2/(2 sigma^2))。
% (a) 在等距中心 c_l=(l-1)/(M-1) 处插值 f*(x)=|x-1/2| 的有限核和，M=5（一致误差 0.061）、M=9（0.024）；
% (b) 中心差分 (K(x,1/2+h)-K(x,1/2))/h（h=0.1、0.02）一致逼近 d/dc K(x,c)|_{c=1/2}=(x-1/2)K(x,1/2)/sigma^2，
%     一致偏差分别为 2.14、0.44。数据由同目录 generate.py 计算。
\begin{tikzpicture}[font=\figurefont\footnotesize]
\pgfplotsset{every axis/.style={width=52mm,height=38mm,axis lines=left,
  x axis line style={draw=FigureStroke,line width=.6pt},y axis line style={draw=FigureStroke,line width=.6pt},
  tick style={draw=FigureStroke},tick label style={font=\figurefont\scriptsize,text=FigureInk},
  xmin=0,xmax=1,xtick={0,0.5,1},ymajorgrids,grid style={draw=FigureGrid,line width=.4pt},
  title style={font=\figurefont\footnotesize},xlabel={$x$},label style={font=\figurefont\scriptsize},
  xlabel style={at={(axis description cs:1,0)},anchor=north west,yshift=1mm},
  legend style={draw=none,fill=none,font=\figurefont\scriptsize,cells={anchor=west},at={(0.5,-0.22)},anchor=north,legend columns=3}}}
\begin{axis}[name=A,ymin=-0.05,ymax=0.55,ytick={0,0.25,0.5},title={(a) 有限核和逼近 $f_*$}]
\addplot[draw=FigureInk,line width=1.1pt] coordinates {(0,0.5) (0.5,0) (1,0.5)};
\addplot[draw=FigureGreen,line width=1.1pt,dashed] table {figures/02-abstract-mathematics/03-function-spaces/tikz/c6-gaussian-sections/gauss-M5.dat};
\addplot[draw=FigureBlue,line width=1.1pt] table {figures/02-abstract-mathematics/03-function-spaces/tikz/c6-gaussian-sections/gauss-M9.dat};
\legend{$f_*$,$M=5$,$M=9$}
\end{axis}
\begin{axis}[at={(A.east)},anchor=west,xshift=12mm,ymin=-5,ymax=5,ytick={-4,0,4},title={(b) 中心差分}]
\addplot[draw=FigureInk,line width=1.1pt,domain=0:1,samples=150] {(x-0.5)/0.0225*exp(-(x-0.5)^2/0.045)};
\addplot[draw=FigureYellow,line width=1.1pt,dash dot] table {figures/02-abstract-mathematics/03-function-spaces/tikz/c6-gaussian-sections/diff-h0.1.dat};
\addplot[draw=FigureBlue,line width=1.1pt,dashed] table {figures/02-abstract-mathematics/03-function-spaces/tikz/c6-gaussian-sections/diff-h0.02.dat};
\legend{导数,$h=0.1$,$h=0.02$}
\end{axis}
\end{tikzpicture}
